\begin{deletedblock}
We converted the two indicators into numerical labels by clustering, as stated in Section~\ref{sec:train-evaluation}.
These labels are referred to as the GL, representing the clustering results and capturing multiple activity levels.
We used k-means as a clustering method~\cite{krishna1999genetic}.
We set $k=3$.
We assigned clusters to which participants belonged as an activity-level indicator.

We also extended the point-cloud dataset, which integrated prolonged video viewing and test-taking, using a participant-similarity-based approach~\cite{shinkuma2019relational} similar to that described in Section~\ref{sec:train-evaluation}.
We extracted two datasets from the obtained six datasets.
The datasets were selected as unique combinations without duplication.
We generated new features from point-cloud features and activity-level measures of the extracted data.
Each point-cloud feature was converted to the absolute value of the difference between the feature values.
The activity levels were converted into binary labels for agreement and disagreement, referred to as ML.
The ML were set to 1 for agreement and 0 for disagreement.
We obtained 15 datasets.
\end{deletedblock}
